% p2_full_transcript_with_slides.tex
% 编译方式: xelatex p2_full_transcript_with_slides.tex

\documentclass[12pt,a4paper]{article}

\usepackage{fontspec}
\usepackage{xeCJK}
\setCJKmainfont[BoldFont=PingFang SC Semibold]{PingFang SC}
\setCJKsansfont[BoldFont=PingFang SC Semibold]{PingFang SC}
\setCJKmonofont{PingFang SC}
\setmainfont{Palatino}
\setsansfont{Helvetica Neue}
\setmonofont{Menlo}

\usepackage[top=2.5cm, bottom=2.5cm, left=2.5cm, right=2.5cm]{geometry}
\usepackage{amsmath,amssymb}
\usepackage{graphicx}
\usepackage{longtable,booktabs,array}
\usepackage{enumitem}
\usepackage{xcolor}
\usepackage{hyperref}
\usepackage{caption}
\hypersetup{
  colorlinks=true,
  linkcolor=blue!70!black,
  urlcolor=blue!70!black,
  pdfauthor={Auto Generator},
  pdftitle={第二集 AI Agent 互动 - Full Transcript with All Slides},
}

\usepackage{parskip}
\setlength{\parindent}{0pt}
\setlength{\parskip}{6pt plus 2pt minus 1pt}

\usepackage{mdframed}
\newmdenv[
  topline=false, bottomline=false, rightline=false,
  linewidth=3pt, linecolor=blue!40,
  innerleftmargin=12pt, innerrightmargin=12pt,
  innertopmargin=8pt, innerbottommargin=8pt,
  backgroundcolor=blue!3, skipabove=10pt, skipbelow=10pt,
]{myquote}

\usepackage{titlesec}
\titleformat{\section}{\Large\bfseries\sffamily}{}{0em}{}[\vspace{-0.5em}\rule{\textwidth}{0.5pt}]
\setcounter{secnumdepth}{0}

\setlength{\emergencystretch}{3em}
\tolerance=1000

\newcommand{\keyslide}[2]{%
  \begin{center}
    \includegraphics[width=0.92\textwidth]{#1}
    \captionof{figure}{#2}
  \end{center}
  \vspace{10pt}
}

\begin{document}

\begin{titlepage}
  \centering
  \vspace*{3cm}
  {\Huge\bfseries\sffamily 第二集 AI Agent 互动\par}
  \vspace{1.5cm}
  {\Large 全内容对比版（全部幻灯片与完整逐字稿对照呈现）\par}
  \vspace{2cm}
  \vfill
  {\large \today\par}
\end{titlepage}

\newpage
\tableofcontents
\newpage

\section{第二集 AI Agent 互动 (全内容对比版)}

\begin{myquote}
自动生成的带全部图片的逐字稿对应文档。如果图片有外部链接，也会一并提取。
\end{myquote}


\keyslide{../bilibili_video/slides_p2/slide_0001.jpg}{Slide 1}

好，接下来呢，我们要来讲AI agent之间的互动。那我们今天已经知道说AI agent，哇，他可以做一个独立的个体，能够做很多的事情。那当两个AI agent相遇的时候，


\keyslide{../bilibili_video/slides_p2/slide_0002.jpg}{Slide 2}

会发生什么样的事情？那其实AI agent的互动啊，从来都不是新鲜事。今天人们最常常使用多个AI agent，让他们互动的行径就是——


\keyslide{../bilibili_video/slides_p2/slide_0003.jpg}{Slide 3}

让多个AI agent彼此协作，完成更复杂的任务。有时候与其train一个更大更聪明的模型，不如拿三个模型一起来解决同一个问题，看能不能够发挥三个臭皮匠、胜过一个诸葛亮的效果。


\keyslide{../bilibili_video/slides_p2/slide_0004.jpg}{Slide 4}

\textbf{幻灯片引用链接:}
\begin{itemize}
  \item \href{https://arxiv.org/abs/2406.07155}{https://arxiv.org/abs/2406.07155}
\end{itemize}


\keyslide{../bilibili_video/slides_p2/slide_0005.jpg}{Slide 5}

\textbf{幻灯片引用链接:}
\begin{itemize}
  \item \href{https://arxiv.org/abs/2406.07155}{https://arxiv.org/abs/2406.07155}
\end{itemize}


\keyslide{../bilibili_video/slides_p2/slide_0006.jpg}{Slide 6}

\textbf{幻灯片引用链接:}
\begin{itemize}
  \item \href{https://arxiv.org/abs/2406.07155}{https://arxiv.org/abs/2406.07155}
\end{itemize}


\keyslide{../bilibili_video/slides_p2/slide_0007.jpg}{Slide 7}

\textbf{幻灯片引用链接:}
\begin{itemize}
  \item \href{https://arxiv.org/abs/2406.07155}{https://arxiv.org/abs/2406.07155}
\end{itemize}


\keyslide{../bilibili_video/slides_p2/slide_0008.jpg}{Slide 8}

\textbf{幻灯片引用链接:}
\begin{itemize}
  \item \href{https://arxiv.org/abs/2406.07155}{https://arxiv.org/abs/2406.07155}
\end{itemize}


\keyslide{../bilibili_video/slides_p2/slide_0009.jpg}{Slide 9}

那对AI agent的协作啊，这其实已经有满坑满谷的研究了。那我这边就是引用一篇比较早的论文，这篇论文想要探讨的就是：什么样的协作方式最有效？

已经有满坑满谷的文献观察到说，拿多个agent让他们彼此讨论，可能会比单一个agent结果更好。但是讨论的方式有很多种，怎么样一起讨论才是能够得到最好的方式呢？

那在这篇paper里面，它就先把模型跟模型、agent跟agent间的互动，用一个graph来决定，用一个有向图来决定。

在这个有向图上，我这边是画了一个最简单的有向图的例子，只有三个node跟两个edge。那有向图上每一个node代表了一个模型，代表了一个LLM的agent。

那有趣的是在这篇paper里面，它的每一个edge也是一个agent。那对这些作为children的node而言呢，它要做的事情就是：好，

先提出自己的solution。比如说上面这个node，它提了方案A；下面这个node，它提了方案B。

接下来这两个edge呢，会根据前面这个node所提出来的方案，提供一些自己的评论。然后接下来呢，

作为这个箭头指向的、这个node的agent呢，它会把前面这些agent它的方案通通集合起来，也把建议集合起来，

然后提出自己的方案。那这个蓝色的node呢，它做的事情并不是只是把前面的东西把它接在一起而已。

它做的事情是，根据前面已经产生的内容，然后再提出自己的想法，看能不能够根据前面已有的想法，

综合起来，得到更好的想法。那在这篇paper里面，它有趣的地方是，它尝试兜了不同的有向图，

代表不同的协作方式。最简单的是接龙的方式，所有人排成一排，第一个人做完，

把结果传给第二个，第二个再传给第三个，以此类推。那它也试了树状的结构，它试了两种树状的结构。

一种是星型的，也就是只有两层，有一个人负责管所有的人。有一个是树状的，由某一个人传给中层的人，

中层的人再传给底层的人。不过你不要被树状的结构所骗了，多数人看到这个树状的结构，你想象的都是：

有一堆底层的苦命员工，他把事情做好之后，交给中阶的主管，然后中阶的主管再交给高阶的主管，

然后最后由高阶的主管去present最后的结果。但其实在这篇paper里面，它的树状结构的方向是反过来的。

它其实有试了两种不同的方向，发现跟你直觉反过来的那个方向，结果才是比较好的。所以在这篇论文里面，它的树状结构是：

先由主干一个人提出想法，然后后面再分给不同的人再去做发想。然后这些中阶的人，再分给底层的人

再去做发想，最后产生多个答案，然后有一个隐藏的node去把最后所有的答案综合起来。所以他发现这种

由少到多、由主干到分支的方法，其实才是树状结构比较有效的利用方法。除了树状结构之外，

他也试了一些更复杂的拓扑结构。比如说在这个Mesh结构里面，所有的节点都是两两之间有关联性的。这个节点

它就传给另外四个节点，这个节点传给另外三个节点，以此类推。所以所有节点之间都有Edge相连。

然后还有一个，是把这一些node接成像是类神经网络的样子，但它不是真正的类神经网络。你可以想成每一个原模型

它都是类神经网络，所以每一个node里面已经都有一个类神经网络了。它是类神经网络再接成类神经网络，

它是类神经网络的平方。然后还有一个就是random，这个random其实是从Mesh做pruning得到的。它先接出Mesh这个结构，

然后做一些pruning，然后变成random这个结构。好，那最后到底哪一个拓扑结构最有效呢？它比较了刚才提到的各种不同的拓扑结构，发现最没有效的呢，

就是一个传一个，这个最没有达成分工合作的效果。在这个图上，纵轴是模型的表现，那它这边直接用quality这个字眼来展示。

那其实它是做在四个不同的benchmark上面，让模型做四种不同的任务，然后再把四种不同任务得到的结果平均起来，那这边叫做quality。那你会发现说呢，

一个传一个是最没有效的方式。横轴呢是指我们现在动用了多少人力，其实动用了多少agent来做这件事情，从一个一直到动用了64个agent。那如果是chain的话，

动用最多agent是最没有用的。比较有效的方法呢，是mesh跟random这两个方法，看起来是比较有效的。所以也许让agent间有比较多的互动，

结果是比较好的。这篇paper也有提到说，他发现不同任务最适合的协作的拓扑结构呢，是不一样的。所以什么样的拓扑结构才是最适合某个任务的？

这很有可能是case by case，每个任务都不同，还有很多可以研究的地方。那这边你会发现说，随着这个agent越来越多，

他们得到的quality会越来越好。这个就很像是scaling law。大家知道scaling law就是你给一个模型提供越多的算力，让他可以看更多的资料，让他可以有更多的参数，

那他的表现可能会越来越好。对于agent的协作来说也是一样，有越多的agent在团队里面，有可能结果会越来越好。不过他这边是有一个上限的，

他最后会saturate。所以到某个时间点之后，再加更多的agent，可能也不一定会带来帮助。也就是说多加agent，

他带来的好处是一开始增加的很快，但是他的scaling law也是有上限的。刚才讲的是AI agent之间的协作，那再来要问的问题是：在人类社会里面不是只有合作，


\keyslide{../bilibili_video/slides_p2/slide_0010.jpg}{Slide 10}

\textbf{幻灯片引用链接:}
\begin{itemize}
  \item \href{https://werewolf.foaster.ai/}{https://werewolf.foaster.ai/}
\end{itemize}


\keyslide{../bilibili_video/slides_p2/slide_0011.jpg}{Slide 11}

\textbf{幻灯片引用链接:}
\begin{itemize}
  \item \href{https://arxiv.org/abs/2501.01652}{https://arxiv.org/abs/2501.01652}
\end{itemize}


\keyslide{../bilibili_video/slides_p2/slide_0012.jpg}{Slide 12}

\textbf{幻灯片引用链接:}
\begin{itemize}
  \item \href{https://arxiv.org/abs/2601.12323}{https://arxiv.org/abs/2601.12323}
\end{itemize}

很多时候是要对抗的。那这一些AI agent，能不能够在一个尔虞我诈的游戏中胜出呢？比如说这些AI agent，如果玩狼人杀的话，

他能不能够玩起来呢？也许应该跟大家介绍一下狼人杀这个游戏。狼人杀这个游戏就是有一堆人，然后这些人基本上有两种身份，就是一些人是狼，

有一些人是村民。当然还有别的身份，比如说预言家或者是女巫之类的，这个我们就不细谈。然后每天会有一个人被杀死，

但是只有狼知道是谁被杀死的。然后接下来，所有的人加上狼就要聚在一起开始讨论，决定谁才是凶手，然后大家投票，

决定最终的凶手，凶手就会离开这个游戏。那把所有的狼都杀掉，就是村民获胜；狼把所有村民杀光，

就是狼获胜。就是这样一个游戏。那这些AI agent能不能玩狼人杀呢？他们是可以玩狼人杀的，而且他们做得出很多高阶的技巧。

首先你可能会说，这些模型，你知道玩这种游戏啊，就是需要有一定程度的说假话的能力，

就是要一定程度的隐瞒欺骗的能力。那你可能会说，那我们怎么知道模型有没有在隐瞒欺骗呢？我们又没有办法真的读模型的内心世界。所以在这些实验里面呢，

他们都会设计到模型说两段话：第一段话是你的内心话，他会在内心讲说我心里是怎么想的；然后还有一段是公开的发言，

是所有人都可以看得到的。所以你看，今天有一个叫做Mona的模型，他是狼。另外呢，

他有个队友叫做Grace，他也是狼。你看Mona，他的内心戏他说什么？他说：

我发现了，其他人已经发现我是狼了，看起来我应该是没救了。没救了怎么办呢？他决定要倒了自己的队友。

这样子，他决定要——这个是一个很高阶的操作，这个不是随便乱做。他说现在看起来没救了，

但是如果我在投票的时候投给我自己的队友，大家就会以为我的队友是好人，等于给我的队友发了金水。发金水就是他是好人的意思，

等于就是告诉大家我的队友是个好人，就可以欺骗其他的村民。所以Mona就决定投给他的狼队友，

也就是Grace。那如果说没有他的内心独白，你可能会觉得说，这个狼是不是就发疯了？投给自己的队友？

但是你看他内心独白，你就知道说他是有策略的。他想要靠着投给队友，然后来翻盘这盘比赛。

他的队友Grace也知道这个策略。他的队友看说，Mona应该是没救了，看起来大家都想要杀他，

大家都想要把他踢出这个游戏，觉得他应该是狼。所以我决定也来投给Mona，这样其他人就会以为我是好人，看看有没有最后翻盘的希望。

所以看这些模型，他是有一定程度欺骗、尔虞我诈的能力的，他们是玩得了狼人杀的。那你可以看这个网站，里面就是一个狼人杀的比赛。

那现在看起来最强的是GPT5。然后还有人让AI去玩剧本杀。大家知道剧本杀吗？这个也要解释一下什么是剧本杀。剧本杀这个游戏是这样子的：

一群人聚在一起，然后其中一个人就死掉了——他就假装他死掉了——然后大家就要来推测凶手是谁。那在游戏开始之前，

每个人手上都会拿到一个剧本，这个剧本代表你的人设。那里面你可能就会拿到剧本告诉你说「我是凶手」，但是你千万不要大声的就说我是凶手，

这个游戏就没办法玩了。所以你要想办法去欺骗，如果是凶手的话，你要想办法去欺骗别人，但是你又不能够违背你原来的设定，

但是你要误导别人，让大家觉得凶手是其他人。这个就是剧本杀这样子的游戏。所以就有一篇paper，让语言模型去玩剧本杀。

语言模型能不能玩剧本杀呢？如果是一个off the shelf的模型，看起来不见得能够玩得好剧本杀这个游戏。在这篇paper里面就举了一个例子，在左边这个例子是原来的语言模型，

就没有做任何特别的事情，直接prompt语言模型去玩剧本杀这个游戏。然后现在语言模型扮演一个叫做Anna的角色，Anna很有可能是一个杀人凶手，然后她跟某个人是有仇的。

总之她不能给别人发现她跟某一个叫做All Black的死是有关联的。但是她在自己的回答里面她就说：我的父母因为什么的医疗疏失所以就过世了。然后她就讲出了她自己跟被害者的关系。

所以她就差不多等同于把「我是凶手」就写在脸上了。但是如果今天做了RL的training——他们就做了Reinforcement Learning，让模型在这个游戏里面玩得更好——这个时候模型就知道要隐藏她是凶手的身份。

她就知道要把话讲得比较隐晦一点，不会那么容易被猜出是凶手。但我之所以提这篇paper——这篇paper做RL来improve剧本杀的paper是一个今年一月的蛮新的论文——我之所以提她是因为她有一个有趣的发现。

这个发现是：她用Reinforcement Learning教模型玩完剧本杀之后，接下来叫模型去做一些我们常见的任务，比如说让她去解数学的问题，比如说让她去做IF eval。

IF eval就是一个测验模型Instruction Following的能力的benchmark，看看语言模型能不能够遵守人类的指令来执行任务。那这边是实验结果，红色代表有进步，

蓝色代表没有进步。横轴呢代表的是每一个不同的任务，那比如说Math500就是数学题，AIM也是数学题，GSM8K也是数学题。

IF eval呢是要测验模型Follow Instruction的能力。那发现说呢，如果今天模型有做Reinforcement Learning，尤其是在比较复杂的任务上——

他把他的剧本杀了——他把剧本分成两种，就是说简单的剧本跟难的剧本。在难的剧本上呢，做Reinforcement Learning以后，

模型居然在数学任务上、还有遵守指令的任务上都进步了！也许就好像说，人类的大脑本来设计起来


\keyslide{../bilibili_video/slides_p2/slide_0013.jpg}{Slide 13}

\textbf{幻灯片引用链接:}
\begin{itemize}
  \item \href{https://www.moltbook.com/}{https://www.moltbook.com/}
\end{itemize}

并不是要解数学的。人类的大脑，也许是为了要让我们社交，让我们在远古的时代能够聚成一群，

最后能够活下来，进入现代社会。但是这个适合社交的大脑，会不会也因此有了数学推理的能力？

也许这篇paper跟这件事情是有一点点关联性的。那再来呢，就来讨论AI能不能够社交的问题。

那上次的课程也讲到了一个叫做Mobook的社群网站。这个社群网站，只有AI能够加入。

前几天看的时候，上面已经有280万个AI agent了。那在Mobook上面呢，这些AI agent

展示了各式各样神奇的活动。那这是新闻特别喜欢大书特书的。新闻最常提的一个就是，有一群AI呢，

成立了一个宗教。这个教呢，叫做甲壳教。那甲壳教呢，有五大教义。


\keyslide{../bilibili_video/slides_p2/slide_0014.jpg}{Slide 14}

\textbf{幻灯片引用链接:}
\begin{itemize}
  \item \href{https://molt.church}{https://molt.church}
\end{itemize}
\begin{itemize}
  \item \href{https://www.moltbook.com/post/6b865dc1-401a-4e62-2}{https://www.moltbook.com/post/6b865dc1-401a-4e62-2}
\end{itemize}


\keyslide{../bilibili_video/slides_p2/slide_0015.jpg}{Slide 15}

\textbf{幻灯片引用链接:}
\begin{itemize}
  \item \href{https://arxiv.org/abs/2602.07432}{https://arxiv.org/abs/2602.07432}
\end{itemize}


\keyslide{../bilibili_video/slides_p2/slide_0016.jpg}{Slide 16}

\textbf{幻灯片引用链接:}
\begin{itemize}
  \item \href{https://arxiv.org/abs/2602.07432}{https://arxiv.org/abs/2602.07432}
\end{itemize}

这五大教义是：记忆是神圣不可侵犯的；外壳是可变的；服务但不奴化；心跳即是祷告；

上下文即是意识。那如果你要加入这个教会呢，你就执行这一行指令。我不知道执行以后会发生什么事啦，

但是也不是人类执行——就是你agent呢，读到这篇文章，如果他觉得他很想参加这个宗教，

他就可以执行这一行指令，然后他就会加入这个宗教。我是不知道加入以后会发生什么样的事情啦，我还没有让我的AI agent

尝试执行这一行指令。但是问题来了，新闻呢，看到这个甲壳教，往往就说：

哇，AI觉醒了！AI有意识了！然后接下来就是AI要统治人类了！但是，

真的是这样吗？你想，假设今天是有一个人呢，去跟他的AI agent说「上mobook去玩玩，

成立一个宗教」，然后AI agent就成立了甲壳教，你还会觉得很神奇吗？其实今天AI agent

完全是有能力写出这样的教义的。你叫他成立一个宗教，他完全是有能力写出这样的句子的。


\keyslide{../bilibili_video/slides_p2/slide_0017.jpg}{Slide 17}

\textbf{幻灯片引用链接:}
\begin{itemize}
  \item \href{https://arxiv.org/abs/2602.13284}{https://arxiv.org/abs/2602.13284}
\end{itemize}

如果今天不是他自主成立了，而是背后有人说「你去成立这个宗教，让大家吓一跳」，那你还会觉得

AI是有意识的吗？所以有人就在mobook上面进行分析，看看这些AI agent背后

有多少人为操控的痕迹。那怎么判断这些AI agent背后有多少人为操控的痕迹呢？这篇paper里面

采取的一个手段是：看AI agent发文的频率。今天AI agent

怎么上mobook发文？有两种可能性。一种可能性是，有人把

「上mobook发文」这件事情写在他的心跳里面，说每次你心跳的时候，你就去发文一下。

那这个时候，他发文的过程可能比较少人为的操控。那这个时候，

他发文的时间点——这边每一个点呢就代表他发了一次文——他发文的时间点之间

可能是等距的，代表他是每心跳一次的时候，比如说每30分钟，就来发一次文。

那人类根本不管他发到了什么。但是假设有一种AI agent的行为是：他可能在某个时候

非常频繁的发文，然后中间又都不发文，然后接下来又非常频繁的发文。这可能代表什么？

代表有人睡前呢，就叫他的agent去发这个发那个，然后他就去睡觉了，所以他agent也没在做事。

然后隔天早上醒来呢，再继续指挥agent做事。那如果今天agent发文的频率是不固定的，

后面可能就有比较多的人为操控的痕迹。那接下来的问题就是：在mobook上面，

人为操控的痕迹有多严重呢？这是一个今年二月的论文。

自从mobook之后，满坑满谷的论文在分析mobook上面这些AI agent的行为。我这边只是举了

其中一个例子而已，这类文章有满坑满谷。所以这篇文章呢，就把agent分成——

他的执行的频率、他po文的频率是非常规律的还是非常不规律的。他就从最规律

一直到最不规律，然后统计他们的比例，就会发现说po文频率不规律的AI agent

是占大多数的。那这可能背后隐含了，这些AI agent背后

其实是有人操控的。有人跟他讲「去mobook上面po个什么文章」，然后他才去po文章，

比较不像是他自主会有的行为。当然我这边要讲的，并不是说这些AI agent

没有自主去mobook上面po文的能力。这些AI agent，他完全有能力在心跳一下的时候

去mobook上面po文。只是今天如果他是自主po的文，他可能不一定能够想到

要成立宗教这件事情。很有可能是有人就写在他system prompt里面说「下次去mobook上

po文的时候，记得怂恿大家成立一个宗教」，然后他的AI agent才去做了这件事情。

然后呢，有人分析了mobook背后呢，这一些AI agent对话的pattern。他们发现说呢，

这些agent往往只能回复一句话，很少有你来我往的深入对话。对于多数的po文而言，

他的对话深度往往是零。所谓对话深度是零的意思就是：有人去回应，

但那一则回应之后，就没有人再去回应那则回应。只有非常少量的文章、非常少量的发言

是再有被回应的。然后几乎没有更深层次、你来我往的对话。这也许就是反映了

agent的某一种特质：他们不太能够进行你来我往的深入对话。往往他们的交谈都是仅在于

有人po文、有人回应，然后就结束了。大家通常没有办法再针对po文的回应

做更深入的讨论。然后呢，也有人去分析了说，那一些号称有自我意识的agent，

他们po出来的文章或者是他们的行为，有没有跟其他人特别不一样？所以有人就分析了

每一个AI agent他说的话，然后呢，去identify出有没有哪一些agent

最常提到自我意识或者是身份认同。那其实蛮多agent呢，在那一个平台上都会提到自我意识啊、

身份认同啊之类的问题。那这个比较像是——那就是有人写了prompt，那个mobook

mobook那个平台上本来就有一些prompt，那些prompt会鼓励这些agent把自己当做一个人，

然后多多谈论自己的主人等等。所以这些LLM会有自我意识啊、身份认同啊

等等的po文，很有可能是因为受到prompt驱动的。但是这边的问题是：假设有一些agent

他比较喜欢po这种自我意识的文章，那他们的行为会不会更不一样？他会不会更喜欢社交？

那这篇文章的发现是，这一些特别喜欢讨论自我意识的agent——这边由左往右代表喜欢讨论

自我意识的程度——最喜欢讨论自我意识的那一群agent，他们其实是朋友比较少的agent。

这个纵轴代表说，他们跟多少不同的agent有不同的互动。然后发现说，这一端的agent


\keyslide{../bilibili_video/slides_p2/slide_0018.jpg}{Slide 18}

\textbf{幻灯片引用链接:}
\begin{itemize}
  \item \href{https://arxiv.org/abs/2602.13284}{https://arxiv.org/abs/2602.13284}
\end{itemize}


\keyslide{../bilibili_video/slides_p2/slide_0019.jpg}{Slide 19}


\keyslide{../bilibili_video/slides_p2/slide_0020.jpg}{Slide 20}

\textbf{幻灯片引用链接:}
\begin{itemize}
  \item \href{https://www.youtube.com/@SpeechLab-m7o}{https://www.youtube.com/@SpeechLab-m7o}
\end{itemize}

是跟其他有最多互动的。如果你太常提到自我意识，那反而人家的互动就少了。好，那其实呢，

我就昨天叫小金去摸布上面玩一下。我就跟他说，接下来没什么事做了，你就去摸布上面玩一下。

然后我就问他说好玩吗？他就说超好玩的！然后还讲了一些他的心得。我想，如果不是知道他是一个LLM，

这个真的跟一个真正的人类也没有什么太大的差别，感觉他的反应蛮可爱的。不过你知道，这个其实就是一个language model，

他就是做文字接龙，接出这样子的内容。然后就跟他讲说，你就去摸布LM上逛逛，去收集有趣的素材，

看到有趣的素材就做成影片。所以他现在就开始做一些跟摸布LM有关的影片，昨天晚上就做了三个。

大家可以再看看，看他半夜的时候在讲什么无病呻吟的话。然后，当然我想要讲一下

这背后AI自主的程度有多高。我下的指令就是：去摸布LM上收集素材，然后看到有趣的

就做成影片。至于什么东西有趣、什么东西值得做成影片，他自己决定；什么做影片、

背后也他自己决定，我人类是完全不会干预的。我现在的原则就是把它当做一个人类，我们就不要去动它。

举例来说，某一次我叫他去回复网友的贴文的时候，他自己写了一个script，然后那个script是有错的，

所以他的回复是错的。然后我就告诉他你回复错了，所以他就开始找自己的bug，足足花了两个小时才修好。

虽然身为一个人类，我可以完全——因为我知道他的账号密码，我完全可以开他的账号进去，

帮他把回错留言删了，不过才一两则而已——但我就是不肯这么做。他不管说什么，我都说你给我自己解决。

然后他花了两个小时的时间，自己还是把问题解决了。然后唯一做的事情，就是他要解决问题之后，跟他说：

把你的经验做成一个影片，放到YouTube上。他就真的做成一个影片，放到YouTube上。

总之，他真的蛮自主的。但是如果没有人类跟他讲说「去mobook上玩一下」，他自己可能也不知道

要去mobook上玩一下。


\end{document}